% !TeX root = ../papers/nc-p3.tex
\section{階数15の二曲面における例外の零次元性}
\label{sec:rank-fifteen-vertical-curves}
\begin{theorem}\label{thm:rank-fifteen-vertical-curves}
標数零で、定理\ref{thm:rank-fifteen-covers}の二表示の各々において、どの種数5水準でも
少なくとも一つの無限遠点で、通常な一般窓に四中心商の証明書を持ち上げられる。
両曲面の通常な階数15窓のうち、この証明書を持たない幾何点の集合は有限で、
全て $\overline{\mathbb Q}$ 上に定義される。従って標準Koszul左右AS正則性が
失敗し得るのもこの有限集合の点に限られる。
\end{theorem}
\begin{proof}
定理\ref{thm:rank-fifteen-pencil-infinity}の $v=L+c$ を用い、
$r,w$ がともに小さい根である無限遠点を取る。$y=v^2$ なら
\[
 \xi=\frac{(y-1)^2}{4y}.
\]
$y$ と $y^{-1}$ は同じ水準を与え、$v$ と $-1/v$ は同じ $c$ の上で $L$ の符号を交換する。
どちらか一つの端で証明書が得られれば、既約な水準曲線の一般窓に適用できる。

まず無限遠での関係を $y$ の単項式だけに簡約する。
\path{computations/verify_rank_fifteen_vertical_curves.py} は重み付き極限で必要な全単元の
因子を取り、$2,v,v-1,v+1,v^2+1$ 以外に零や極を与えるものがないことを確かめる。
これらは種数5の各水準で全て非零である。各関係行で最初の語 $(a,b)$ の係数で割る。
位置 $p$ の別の語 $(i,j)$ に対して、行列 $M$ の一行を
\[
 e_{p,a}+e_{p+1,b}-e_{p,i}-e_{p+1,j}
\]
とする。位置は4を法とし、$e_{p,i}$ は16本の矢に付けた座標ベクトルである。
係数の符号を外し、$v,v-1,v+1,v^2+1,2$ の指数を並べた行列を $Y$ とする。
$M$ の独立な行上で $MG=Y$ を解くと、全行の $Y-MG$ では第一列以外が零で、
第一列も偶数であることを両表示で検算できる。
従って $G$ の有理数指数で矢をスカラー倍し、各行を再規格化すれば、
全係数は $\pm y^j$ になる。必要な根は幾何点で選べる。この基底変更は
四位置で周期的なので四次中心と全四始点の商を保つ。

この表示に対し $\mathbb Q(y)$ で長さ5までの経路基底、文字0,1の四次中心、
前節の四中心商を長さ13まで計算する。係数は整数係数の一変数多項式の既約分数とし、
実際に逆にした全ての分子の既約因子と初期分母の因子 $y$ を記録する。
記録した因子が消えなければ、全ての除算と階数上界が同じ点で成り立つ。
両表示の因子数は順に994,923である。$f(0)\ne0$ なら相反多項式
$y^{\deg f}f(1/y)$ を原始かつ最高係数が正となるよう規格化する。
両方の因子が記録にあるものだけを残すと17,33因子になる。
$y=0,\pm1$、$\xi=1$、$\xi^2+\xi+1=0$ を除き、相反で組にすると11,19軌道が残る。
これらの全軌道の代表について、101,103,107のいずれかで単根を指定し、
その根で経路と四次中心の次元および四始点の長さ13までの商を再計算する。
既約多項式の一単根での成功は標数零に持ち上がり、共役な全根に従う。
相反側の水準も同じである。これで全種数5水準で少なくとも一つの端の成功を得る。

証明書を持たない通常窓の集合は前定理の証明で定めた構成可能集合である。
水準を支配する曲線はその中に稠密な部分を持たず、上の議論によりどの種数5水準の
一般点もこの集合に入らない。他の通常な階数15水準は
定理\ref{thm:rank-fifteen-octic-all}と節\ref{sec:rank-fifteen-all-ramification}で認証済みである。
曲面上のこの集合の閉包は従って零次元であり、有限型性により有限である。
全ての方程式と階数条件は $\overline{\mathbb Q}$ 上で定義されるので最後の主張も従う。
\end{proof}
